% Printed 518; principal PDF 525; Google PDF 529.
\par\clearpage
\noindent\makebox[0pt][l]{\fontsize{13}{16}\selectfont 518}\hfill{\small\textbf{BEMERKUNGEN.}}\hfill\null\par\medskip

Für die reciproken Gewichte erhält man aus den obigen Formeln
\[
\begin{alignedat}{3}
 \text{bei }n&=4\quad&
 \frac{1}{P}&=\frac{1}{rr}\left(\tfrac14g-\tfrac14\frac{3g+2}{2g+1}\right)&
 &=\frac{kk}{rr}\left(\tfrac14+\tfrac14\frac{kk+2}{kk(2kk+5)}\right)\\
 n&=6&
 \frac{1}{P}&=\frac{1}{rr}\left(\tfrac16g-\tfrac16\frac{5gg+4g-3}{3gg+2g-2}\right)&
 &=\frac{kk}{rr}\left(\tfrac16+\tfrac16\frac{k^4+4kk+3}{kk(3k^4+14kk+14)}\right)\\
 n&=3&
 \frac{1}{P}&=\frac{1}{rr}\left(\tfrac13g-\tfrac23\frac{2g+3}{3g+4}\right)&
 &=\frac{kk}{rr}\left(\tfrac13+\tfrac23\frac{kk+3}{kk(3kk+10)}\right)\\
 n&=5&
 \frac{1}{P}&=\frac{1}{rr}\left(\tfrac15g-\tfrac25\frac{4gg+7g}{5gg+8g-1}\right)&
 &=\frac{kk}{rr}\left(\tfrac15+\tfrac25\frac{k^4+5kk+5}{kk(5k^4+28kk+35)}\right).
\end{alignedat}
\]
\textsc{Gauss} gibt dagegen in der vorstehenden Notiz an: »Wahrscheinlicher Fehler des Resultats, den des Ablesens $=1$ gesetzt:
\[
 \surd\bigl((kk+1)c^{(n)}-c^{(n-1)}\bigr)
 =k\left\delimiter"270370\relax\vphantom{\frac{1}{n}}\right.\!
 \left(\frac{1}{n}+\frac{2c^{(n)}}{nkk}\right).
\]
\indent Auch ist, wenn die Coefficienten mit $c^{(n)}$, $c^{(n-1)}$, etc. bezeichnet werden,
\[
 (nkk+(n-2))c^{(n)}-nc^{(n-1)}=kk.\text{«}
\]

Für den Fall, dass am Anfang, nach der $1^{\text{ten}}$, $(n-1)^{\text{ten}}$ und $n^{\text{ten}}$ Beobachtung abgelesen ist, möge eine Ableitung des Winkelwerthes und seines reciproken Gewichts hier folgen.

Es seien\par
\noindent
\begin{tabular}{@{}c@{\quad}p{.80\linewidth}@{}}
 $v_0$, $v_1$, $v_{n-1}$, $v_n$ & die Fehler der Ablesungen $A$, $B$, $Y$, $Z$;\\
 $\delta_1$, $\delta_{n-1}$, $\delta_n$ & die Summe der Fehler der Visuren zwischen den Ablesungen $A$ und $B$, $B$ und $Y$, $Y$ und $Z$;\\
 $x$ & der gesuchte Winkelwerth.
\end{tabular}\par

Das Gewicht eines Fehlers der Ablesung sei wie vorher $1$, das Gewicht eines einzelnen Visurfehlers sei $p$. Die Gewichte von $\delta_1$, $\delta_{n-1}$, $\delta_n$ sind demnach $\dfrac{p}{2}$, $\dfrac{p}{2(n-2)}$, $\dfrac{p}{2}$, da sich $\delta_1$ und $\delta_n$ je aus $2$ Einstellungen, $\delta_{n-1}$ aber aus $n-2$ Einstellungen zusammensetzen. Man setze jetzt $\dfrac{p}{2}=\dfrac{1}{kk}$.

Zwischen den Fehlern bestehen die Gleichungen
\[
\makebox[\linewidth]{\hspace*{\parindent}\makebox[0pt][l]{\small 1)}\hfill
\(\displaystyle
 \begin{aligned}
  \delta_1 &= -A+B-x-v_0+v_1\\
  \delta_{n-1} &= -B+Y-(n-2)x-v_1+v_{n-1}\\
  \delta_n &= -Y+Z-x-v_{n-1}+v_n.
 \end{aligned}
\)\hfill\hspace*{\parindent}}
\]
\indent Es muss
\[
\makebox[\linewidth]{\hspace*{\parindent}\makebox[0pt][l]{\small 2)}\hfill
\(\displaystyle
 v_0v_0+v_1v_1+v_{n-1}v_{n-1}+v_nv_n+
 \frac{1}{kk}\left(\delta_1\delta_1+
 \frac{1}{n-2}\delta_{n-1}\delta_{n-1}+\delta_n\delta_n\right)=\min.
\)\hfill\hspace*{\parindent}}
\]
werden.

Setzt man für den Augenblick
\[
 -A+B=L_1,\qquad -B+Y=L_2,\qquad -Y+Z=L_3,
\]
so lauten die Normalgleichungen, wenn man gleichzeitig mit $kk$ multiplicirt:
{\small
\[
\begin{array}{@{}l@{\;}r@{\;}c@{\;}c@{\;}c@{\;}c@{\;}c@{\;}c@{}}
 & L_1+L_2+L_3 &=& nx & +v_0 & * & * & -v_n\\
 & \multicolumn{1}{l@{\;}}{L_1} &=& +x & +(kk+1)v_0 & -v_1 & * & *\\
 \text{3)} & -L_1+\dfrac{1}{n-2}L_2 &=& * & -v_0 &
  +\left(kk+1+\dfrac{1}{n-2}\right)v_1 & -\dfrac{1}{n-2}v_{n-1} & *\\[3pt]
 & -\dfrac{1}{n-2}L_2+L_3 &=& * & * &
  -\dfrac{1}{n-2}v_1 & +\left(kk+1+\dfrac{1}{n-2}\right)v_{n-1} & -v_n\\[3pt]
 & -L_3 &=& -x & * & * & -v_{n-1} & +(kk+1)v_n.
\end{array}
\]
}

% Printed 519; principal PDF 526; Google PDF 530.
\par\clearpage
\noindent\null\hfill{\small\textbf{REPETITIONSBEOBACHTUNGEN.}}\hfill\makebox[0pt][r]{\fontsize{13}{16}\selectfont 519}\par\medskip

Addirt man die letzten $4$ Gleichungen, so folgt
\[
 v_0+v_1+v_{n-1}+v_n=0;
\]
nach der zweiten und fünften Gleichung ist mithin:
\[
 L_1-L_3=(kk+2)(v_0+v_n)=-(kk+2)(v_1+v_{n-1}).
\]

Die erste Normalgleichung gibt
\[
 \delta_1+\delta_{n-1}+\delta_n=0.
\]

Subtrahirt man die fünfte Gleichung von der zweiten und die vierte von der dritten, so hat man
\[
\begin{array}{@{}l@{\;}r@{\;}r@{\;}c@{\;}c@{\;}r@{\;}l@{}}
 L_1 && +L_3 &=& 2x & +(1+kk)(v_0-v_n) & -\hphantom{\dfrac{kk(n-2)+n}{n-2}}(v_1-v_{n-1})\\
 -L_1 & +\dfrac{2}{n-2}L_2 & -L_3 &=& * & -(v_0-v_n) & +\dfrac{kk(n-2)+n}{n-2}(v_1-v_{n-1}),
\end{array}
\]
woraus sich ergibt:
\[
 (kk(n-2)+2)(L_1+L_3)+2L_2=2(kk(n-2)+n)x+
 (k^2(n-2)+2kk(n-1)+2)(v_0-v_n).
\]

Eliminirt man nun vermittelst dieser Gleichung $v_0-v_n$ aus der ersten der Gleichungen 3), so folgt, wenn man zugleich für die $L$ ihre Werthe einsetzt:
\[
 x=\frac{((n-2)kk+n)(Z-A)+(n-2)(Y-B)}{n(n-2)kk+2nn-4n+4}.
\]

Um das reciproke Gewicht von $x$ zu erhalten, hat man an Stelle der Constanten der Gleichungen 3) (weil die ursprünglichen Normalgleichungen mit $kk$ multiplicirt sind) der Reihe nach $kk$, $0$, $0$, $0$, $0$ zu setzen. Der damit erhaltene Werth von $x$ gibt das reciproke Gewicht an. Man erhält
\[
 \frac{1}{P}=\frac{(n-2)k^4+2(n-1)kk+2}{n(n-2)kk+2nn-4n+4}.
\]

Setzt man dagegen, wie es \textsc{Gauss} in der letzten Formel auf S.~514 und in den beiden folgenden Formeln auf S.~515 gethan hat, das Gewicht der zu einer Beobachtung gehörigen Visurfehler $=1$, das Gewicht eines Ablesefehlers $=kk$, so ist im Nenner des vorstehenden Ausdrucks noch $kk$ als Factor hinzu zu fügen.

Für $n=12$ und $kk=\frac19$ wird alsdann z.~B. das Gewicht des Winkelwerths:
\[
 P=\frac{\frac19\left\{12.10.\frac19+2.144-4.12+4\right\}}{10.\frac{1}{81}+2.11.\frac19+2}
 =6{,}26.
\]

Wäre aber nach jeder Beobachtung abgelesen, so hätte man nach S.~516/517 für $kk=\frac19$ zunächst
\[
 \mu_1=1,\qquad \mu_2=2{,}111,\qquad \mu_3=3{,}457,\qquad
 \mu_4=5{,}187,\qquad \mu_5=7{,}493,\qquad \mu_6=10{,}631;
\]
und weiter mit $n=12$:
\[
 (nkk+n-2)\mu_m-n\mu_{m-1}=30{,}574=kkM.
\]

Damit ergibt sich, da in diesem Falle der Ausdruck für das reciproke Gewicht noch durch $kk$ zu dividiren, und da ausserdem $r=1$ ist:
\[
 \frac{1}{P}=\frac{1}{n}+\frac{2\mu_m}{nkkM}
\]
oder

% Printed 520; principal PDF 527; Google PDF 531.
\par\clearpage
\noindent\makebox[0pt][l]{\fontsize{13}{16}\selectfont 520}\hfill{\small\textbf{BEMERKUNGEN.\quad REPETITIONSBEOBACHTUNGEN.}}\hfill\null\par\medskip
\[
 P=\frac{12kkM}{kkM+2\mu_6}=\frac{12.30{,}574}{51{,}836}=7{,}08.
\]
\indent Vergl. dazu den Schluss des Briefes an \textsc{Bessel}, S.~516.

Ist nur im Anfang und nach der $n^{\text{ten}}$ Beobachtung abgelesen, so findet man, entweder direct wie vorher oder indem man in der allgemeinen Formel auf S.~517 $n=2m+1=1$ setzt, für das reciproke Gewicht von $x=\dfrac{Z-A}{n}$:
\[
 \frac{1}{P}=\frac{nkk+2}{nn}.
\]

Dabei ist das Gewicht eines Ablesungsfehlers $=1$ und das Gewicht des Gesammtfehlers des Pointirens $=\dfrac{1}{nkk}$ gesetzt.

Das \textsc{Svanberg}sche Verfahren (Exposition des opérations faites en Laponie, pour la détermination d'un arc du méridien, en 1801, 1802 et 1803, Stockholm 1805, p.~29/33) besteht darin, dass zur Ableitung des Winkelwerths $x$ aus $(n+1)$ Ablesungen $l_0$, $l_1$, \ldots, $l_n$ alle möglichen Differenzen $l_i-l_k=(i-k)x$, $i>k$, gebildet werden, aus deren Summe man erhält
\[
 \frac{n(n+1)(n+2)}{1.2.3}x=\sum_{i=0}^{n}(2i-n)l_i.
\]

Wie \textsc{Gauss} in dem vorstehenden Briefe an \textsc{Bessel} bemerkt, findet man dasselbe Resultat, wenn man die Pointirungsfehler, $\delta$, gleich Null setzt. In der That ergibt sich dieser Werth von $x$ auch mit Hülfe der Fehlergleichungen:
\[
 l_0+v_0=u,\quad l_1+v_1=u+x,\quad l_2+v_2=u+2x,\quad \ldots,\quad l_n+v_n=u+nx,
\]
worin $u$ die Entfernung des Anfangspunktes der Theilung von der ersten Ablesestelle bedeutet.\par
\hfill\textsc{Krüger.}\hspace*{\parindent}\par

\vspace{2.5cm}
\begin{center}\rule{4cm}{.4pt}\end{center}

% Unfoliated [521]; principal PDF 528; Google PDF 532.
\par\clearpage
\vspace*{.11\textheight}
\begin{center}
 {\Large BEMERKUNGEN ZUM NEUNTEN BANDE.}\par
 \vspace{1em}
 \rule{1.5cm}{.4pt}
\end{center}
\vspace{1em}

Der vorliegende neunte Band von \textsc{Gauss}' Werken ist von Herrn L.~\textsc{Krüger} in Potsdam bearbeitet worden, welcher von Herrn A.~\textsc{Börsch} bei der Sichtung des Materials unterstützt wurde; er enthält einestheils den Abdruck derjenigen von \textsc{Gauss} publicirten geodätischen Schriften, welche in Band IV nicht Platz gefunden haben, anderntheils eine Bearbeitung des geodätischen Nachlasses.

Die von \textsc{Gauss} hinterlassenen geodätischen Notizen beziehen sich fast ausschliesslich auf die \textit{hannoversche Gradmessung} und auf Fragen, die durch sie und auch durch ihre Vorgängerin, die dänische Gradmessung, veranlasst sind. Herr L.~\textsc{Krüger} hat sich der mühsamen Arbeit unterzogen, den Inhalt der im Nachlass vorhandenen zusammenhangslosen Blätter, auf denen \textsc{Gauss} die \textit{Projectionsmethode der hannoverschen Gradmessung} entwickelt hat, ebenso wie die auf die \textit{Triangulirung} selbst bezüglichen Aufzeichnungen zu einem geordneten Ganzen zusammen zu setzen. Der \textit{Fortgang der Arbeiten zur Gradmessung und zu der sich anschliessenden Landesvermessung} ist durch eine Zusammenstellung aller wichtigen sie betreffenden Veröffentlichungen, officiellen Berichte und Briefstellen ausführlich dargestellt. Im Einzelnen sind die den verschiedenen Notizen zugefügten Bemerkungen des Bearbeiters nachzusehen; eine Übersicht der wesentlichsten Theile dieses Bandes enthält auch der vom Unterzeichneten im April 1902 erstattete »Bericht über den Stand der Herausgabe von \textsc{Gauss}' Werken« (Nachrichten der K.~Gesellschaft der Wissenschaften zu Göttingen. Geschäftliche Mittheilungen 1902, Heft 1).

Die allgemeine Redaction lag, wie bei Band VIII, in den Händen von Herrn M.~\textsc{Brendel}. Der Abdruck der Nachlassstellen schliesst sich so eng wie möglich an die Originale an; bezüglich der Orthographie wurde, wie in den frühern Bänden, diejenige Schreibweise gewählt, die \textsc{Gauss} am häufigsten gebraucht zu haben schien, und diese durch den ganzen Band beibehalten. Ebenso wie früher sind Einschaltungen des Bearbeiters in den Text der Originale durch eckige Klammern [] kenntlich gemacht oder in den Bemerkungen besonders erwähnt, sowie Abdrücke von Notizen und Briefstellen, die nicht von \textsc{Gauss} herrühren, in geschweifte Klammern \{\} gesetzt.

Die Briefstellen wurden nach den Originalen abgedruckt, soweit diese vorhanden sind. Der Nachweis der aus dem Nachlass zum Abdruck gelangten Notizen findet sich, wie andere redactionelle Mittheilungen, in den Bemerkungen zu den einzelnen Abtheilungen.\par
\hfill F.~\textsc{Klein.}\hspace*{\parindent}\par

\begin{center}\rule{4cm}{.4pt}\end{center}
\noindent\makebox[0pt][l]{\small IX.}\makebox[.86\linewidth][r]{66}\par

% Unfoliated [522]; principal PDF 529; Google PDF 533.
\par\clearpage
\begingroup
\small
\newcount\IXContentsDot
\newcommand{\IXContentsEntry}[5]{%
  \par\noindent\hbox to \linewidth{%
    \hspace*{#1}\strut#2%
    \ifnum#3>0\relax
      \hskip.35em
      \IXContentsDot=0\relax
      \loop\ifnum\IXContentsDot<#3\relax
        \hfil.\advance\IXContentsDot by 1\relax
      \repeat
      \hskip.55em
    \fi
    \hfil\makebox[2.35em][c]{#4}\makebox[1.8em][r]{#5}%
  }\par
}
\newcommand{\IXContentsText}[2]{%
  \par\noindent\hspace*{#1}\strut#2\par
}
\newcommand{\IXContentsHeading}[1]{%
  \par\medskip\noindent\makebox[\linewidth][c]{#1}\par\smallskip
}

\vspace*{.11\textheight}
\begin{center}
 {\Large\textls[200]{INHALT.}}\par
 \vspace{1em}
 {\large GAUSS WERKE BAND IX.\\
 GEODÄSIE.\quad FORTSETZUNG VON BAND IV.}\par
 \vspace{1em}
 \rule{1.5cm}{.4pt}
\end{center}
\vspace{1em}
\begin{center}
 BESTIMMUNG DES BREITENUNTERSCHIEDES ZWISCHEN DEN STERNWARTEN\\
 VON GÖTTINGEN UND ALTONA DURCH BEOBACHTUNGEN\\
 AM RAMSDENSCHEN ZENITHSECTOR.
\end{center}
\IXContentsEntry{1.8em}{Einleitung}{29}{Seite}{5}
\IXContentsEntry{1.8em}{Die beobachteten Sterne}{24}{—}{8}
\IXContentsEntry{1.8em}{Die Beobachtungen}{26}{—}{10}
\IXContentsText{1.8em}{Resultate}
\IXContentsText{3.6em}{Einfachste Combination der Beobachtungen zur Bestimmung des Breitenunterschiedes,}
\IXContentsEntry{5.4em}{Art.~1}{28}{—}{29}
\IXContentsEntry{3.6em}{Genauigkeit der Beobachtungen, Art.~2}{17}{—}{30}
\IXContentsEntry{3.6em}{Collimationsfehler, Art.~3}{23}{—}{31}
\IXContentsEntry{3.6em}{Absolut vortheilhafteste Combination der Beobachtungen, Art.~4—7}{7}{—}{32}
\IXContentsEntry{3.6em}{Berücksichtigung der unregelmässigen Theilungsfehler, Art.~8—11}{8}{—}{36}
\IXContentsEntry{3.6em}{Lage der Beobachtungsplätze, Art.~12}{18}{—}{40}
\IXContentsText{3.6em}{Bestimmung der absoluten Polhöhe der Göttinger Sternwarte aus Beobachtungen des}
\IXContentsEntry{5.4em}{Nordsterns am \textsc{Reichenbach}schen Meridiankreise, Art.~13—17}{7}{—}{40}
\IXContentsEntry{3.6em}{Endresultat der hannoverschen Gradmessung, Art.~18—20}{11}{—}{47}
\IXContentsText{3.6em}{Vergleichung der Declinationen der beobachteten Zenithalsterne mit \textsc{Bradley}s und}
\IXContentsEntry{5.4em}{\textsc{Piazzi}s Bestimmungen, Art.~21}{19}{—}{50}
\IXContentsEntry{1.8em}{Breitenbestimmung der Sternwarte Seeberg}{18}{—}{52}
\IXContentsEntry{1.8em}{Zusatz zu Art.~20, S.~48}{24}{—}{56}
\IXContentsText{0pt}{\textit{Anzeige.}}
\IXContentsText{1.8em}{Bestimmung des Breitenunterschiedes zwischen den Sternwarten von Göttingen und Altona}
\IXContentsEntry{3.6em}{durch Beobachtungen am \textsc{Ramsden}schen Zenithsector}{12}{—}{59}
\IXContentsEntry{1.8em}{Bemerkungen}{28}{—}{63}

% Printed 523; principal PDF 530; Google PDF 534.
\par\clearpage
\noindent\null\hfill{\small\textbf{INHALT.}}\hfill\makebox[0pt][r]{\fontsize{13}{16}\selectfont 523}\par\medskip

\IXContentsHeading{ERDELLIPSOID UND GEODÄTISCHE LINIE.}
\IXContentsText{0pt}{\textit{Nachlass.}}
\IXContentsEntry{1.8em}{Das Erdellipsoid}{27}{Seite}{67}
\IXContentsEntry{1.8em}{Gleichung der Verticalebene des Rotationsellipsoids}{15}{—}{70}
\IXContentsEntry{1.8em}{Gleichung des Rotationsellipsoids in Beziehung auf eine berührende Ebene}{6}{—}{70}
\IXContentsEntry{1.8em}{Bemerkungen}{28}{—}{71}
\IXContentsEntry{1.8em}{Begründung meiner Theorie der geodätischen Linie}{15}{—}{72}
\IXContentsEntry{1.8em}{Kürzeste Linie auf dem Sphäroid}{21}{—}{74}
\IXContentsEntry{1.8em}{Geodätische Linie}{26}{—}{78}
\IXContentsEntry{1.8em}{Geodätische Übertragung von Breite, Länge und Azimuth}{12}{—}{80}
\IXContentsEntry{1.8em}{Geodätische Übertragung auf der Kugel}{19}{—}{88}
\IXContentsText{1.8em}{Berechnung der linearen Länge der geodätischen Linie und ihrer Azimuthe aus den geo-}
\IXContentsEntry{3.6em}{graphischen Coordinaten}{23}{—}{89}
\IXContentsEntry{1.8em}{Vollkommen genaue Formeln für ein Dreieck auf dem elliptischen Sphäroid}{6}{—}{92}
\IXContentsText{1.8em}{Übertragung der geographischen Lage vermittelst der Sehne und des Azimuths des Vertical-}
\IXContentsEntry{3.6em}{schnittes}{28}{—}{93}
\IXContentsEntry{1.8em}{Der Unterschied zwischen dem geodätischen und dem beobachteten Azimuth}{6}{—}{94}
\IXContentsEntry{1.8em}{Reduction des astronomischen Azimuthes auf das geodätische}{11}{—}{95}
\IXContentsEntry{1.8em}{Bemerkungen}{28}{—}{96}
\IXContentsText{0pt}{\textit{Briefwechsel.}}
\IXContentsEntry{1.8em}{Änderung der Polhöhe mit der Höhe: \textsc{Gauss} an J.~J. \textsc{Baeyer} 1853 Juni 22}{6}{—}{99}
\IXContentsEntry{1.8em}{Bemerkungen}{28}{—}{102}
\IXContentsText{0pt}{\textit{Nachlass.}}
\IXContentsEntry{1.8em}{Reduction der sphärischen Dreieckswinkel $A$, $B$, $C$ auf die Chordenwinkel $\mathfrak{A}$, $\mathfrak{B}$, $\mathfrak{C}$}{3}{—}{103}
\IXContentsText{1.8em}{Bedingung dafür, dass 3 Punkte auf der Oberfläche einer Kugel auf einem grössten Kreise}
\IXContentsEntry{3.6em}{liegen}{29}{—}{103}
\IXContentsEntry{1.8em}{Bemerkungen}{28}{—}{104}

\IXContentsHeading{CONFORME DOPPELPROJECTION DES SPHÄROIDS AUF DIE KUGEL UND DIE EBENE.}
\IXContentsText{0pt}{\textit{Nachlass.}}
\IXContentsEntry{1.8em}{Das elliptische Sphäroid auf die Kugel übertragen}{15}{—}{107}
\IXContentsEntry{1.8em}{Bemerkungen}{28}{—}{115}
\IXContentsEntry{1.8em}{Stereographische Projection der Kugel auf die Ebene}{14}{—}{117}
\IXContentsEntry{1.8em}{Bemerkungen}{28}{—}{121}
\IXContentsEntry{1.8em}{Übertragung der Kugel auf die Ebene durch \textsc{Mercator}s Projection}{8}{—}{123}
\IXContentsEntry{1.8em}{Bemerkungen}{28}{—}{132}
\medskip
\IXContentsEntry{1.8em}{Stereographische Darstellung des Sphäroids in der Ebene}{12}{—}{133}
\IXContentsEntry{1.8em}{Bemerkungen}{28}{—}{134}

\vspace{1.7cm}
\noindent\makebox[.86\linewidth][r]{66\textsuperscript{*}}\par

% Printed 524; principal PDF 531; Google PDF 535.
\par\clearpage
\noindent\makebox[0pt][l]{\fontsize{13}{16}\selectfont 524}\hfill{\small\textbf{INHALT.}}\hfill\null\par\medskip

\IXContentsHeading{CONFORME ÜBERTRAGUNG DES SPHÄROIDS AUF DEN KEGELMANTEL.}
\IXContentsText{0pt}{\textit{Nachlass.}}
\IXContentsEntry{1.8em}{Zur zweiten Darstellungsart des Sphäroids, auf einen Parallelkreis bezogen}{6}{Seite}{137}
\IXContentsEntry{1.8em}{Bemerkungen}{28}{—}{140}

\IXContentsHeading{CONFORME ABBILDUNG DES SPHÄROIDS IN DER EBENE}
\noindent\makebox[\linewidth][c]{(PROJECTIONSMETHODE DER HANNOVERSCHEN LANDESVERMESSUNG).}\par\smallskip
\IXContentsText{0pt}{\textit{Nachlass.}}
\IXContentsEntry{1.8em}{Berechnung der geographischen Breite und Länge aus den ebenen rechtwinkligen Coordinaten}{0}{—}{143}
\IXContentsEntry{1.8em}{Berechnung der Meridianconvergenz aus den ebenen rechtwinkligen Coordinaten}{4}{—}{146}
\IXContentsEntry{1.8em}{Formeln zur numerischen Berechnung der Länge, Breite und Meridianconvergenz}{4}{—}{148}
\IXContentsEntry{1.8em}{Berechnung des Vergrösserungsverhältnisses $n$}{17}{—}{152}
\IXContentsEntry{1.8em}{Beziehungen zwischen $x$, $y$ und $\xi$, $\lambda$}{20}{—}{155}
\IXContentsEntry{1.8em}{Berechnung der ebenen rechtwinkligen Coordinaten aus der geographischen Breite und Länge}{0}{—}{156}
\IXContentsEntry{1.8em}{Berechnung der Meridianconvergenz aus den geographischen Coordinaten}{7}{—}{158}
\IXContentsEntry{1.8em}{Die Reduction des Azimuths auf dem Sphäroid auf das Azimuth in plano}{6}{—}{159}
\IXContentsText{1.8em}{Der Unterschied zwischen der Projection der geodätischen Linie und der ihre Endpunkte}
\IXContentsText{3.6em}{verbindenden Geraden bei der conformen Darstellung einer krummen Fläche in der}
\IXContentsEntry{3.6em}{Ebene}{29}{—}{162}
\IXContentsEntry{1.8em}{Zur Transformation der Coordinaten}{20}{—}{168}
\IXContentsEntry{1.8em}{Reihen zwischen $\varphi$, $\psi$ und $\omega$}{22}{—}{171}
\IXContentsEntry{1.8em}{Zur Berechnung von $\log\cos\varphi$}{22}{—}{180}
\IXContentsEntry{1.8em}{Berechnung von $\log(1-ee\sin\varphi^2)$, $\log\frac{(1-ee\sin\varphi^2)^2}{aa(1-ee)}$ und $\frac{a\cos\varphi}{\surd(1-ee\sin\varphi^2)}$}{7}{—}{181}
\IXContentsEntry{1.8em}{Numerische Werthe der Coefficienten in den Reihen zwischen $\varphi$, $\psi$ und $\omega$}{6}{—}{182}
\IXContentsText{1.8em}{Berechnung der ebenen rechtwinkligen Coordinaten aus den geographischen Coordinaten mit}
\IXContentsEntry{3.6em}{Hülfe der Reihen zwischen $\varphi$, $\psi$ und $\omega$}{17}{—}{185}
\IXContentsEntry{1.8em}{Berechnung der Länge und Breite aus den ebenen Coordinaten}{10}{—}{191}
\IXContentsEntry{1.8em}{Die Darstellung der Oberfläche des Sphäroids in der Ebene}{11}{—}{193}
\IXContentsEntry{1.8em}{Bemerkungen}{28}{—}{195}
\IXContentsText{0pt}{\textit{Briefwechsel.}}
\IXContentsText{1.8em}{Über die Formeln für die hannoversche Landesvermessung:}
\IXContentsEntry{.36\linewidth}{\textsc{Gauss} an \textsc{Schumacher} 1830 April 18}{7}{—}{205}
\IXContentsEntry{.36\linewidth}{\textsc{Gauss} an \textsc{Schumacher} 1830 April 30}{7}{—}{212}
\IXContentsEntry{.36\linewidth}{\textsc{Gauss} an \textsc{Schumacher} 1831 Mai 17}{7}{—}{213}
\IXContentsEntry{.36\linewidth}{\textsc{Gauss} an \textsc{Schumacher} 1831 Juni 25}{7}{—}{215}
\IXContentsEntry{.36\linewidth}{\textsc{Gauss} an \textsc{Schumacher} 1838 Dec.~9}{7}{—}{217}
\IXContentsEntry{1.8em}{Bemerkungen}{28}{—}{218}

\IXContentsHeading{TRIGONOMETRISCHE PUNKTBESTIMMUNG.}
\IXContentsText{0pt}{\textit{Nachlass.}}
\IXContentsText{1.8em}{Endresultat für den Ort eines Punktes in einer Ebene, der von drei bekannten aus ange-}
\IXContentsEntry{3.6em}{schnitten ist}{27}{—}{221}

% Printed 525; principal PDF 532; Google PDF 536.
\par\clearpage
\noindent\null\hfill{\small\textbf{INHALT.}}\hfill\makebox[0pt][r]{\fontsize{13}{16}\selectfont 525}\par\medskip

\IXContentsText{1.8em}{Bestimmung der Lage eines Punktes $P^\circ$ aus der Lage dreier anderer: $P$, $P'$, $P''$, wo jener}
\IXContentsEntry{3.6em}{beobachtet}{27}{Seite}{223}
\IXContentsEntry{1.8em}{Ausgleichung dreier Schnitte}{23}{—}{224}
\IXContentsEntry{1.8em}{Zur Ausgleichung dreier Schnitte}{21}{—}{226}
\IXContentsText{1.8em}{Bestimmung eines Nebenpunktes (Scheessel) aus den Beobachtungen auf Hauptdreieckspunkten}
\IXContentsEntry{3.6em}{(Litberg, Wilsede, Bottel, Bullerberg und Brüttendorf)}{12}{—}{228}
\IXContentsText{0pt}{\textit{Abhandlung.}}
\IXContentsEntry{1.8em}{Anwendung der Wahrscheinlichkeitsrechnung auf eine Aufgabe der praktischen Geometrie}{1}{—}{231}
\IXContentsText{0pt}{\textit{Nachlass.}}
\IXContentsText{1.8em}{Bestimmung der Lage des Punktes $X$ durch Beobachtung der Winkel $\alpha$ und $\beta$ zwischen}
\IXContentsEntry{3.6em}{3 gegebenen Punkten $A$, $B$, $C$}{21}{—}{238}
\IXContentsEntry{1.8em}{Orientirung des Messtisches}{23}{—}{239}
\IXContentsEntry{1.8em}{Aufgabe der praktischen Geometrie}{20}{—}{239}
\IXContentsEntry{1.8em}{Bemerkungen}{28}{—}{240}

\IXContentsHeading{AUSGLEICHUNG EINFACHER FIGUREN.}
\IXContentsText{0pt}{\textit{Nachlass.}}
\IXContentsEntry{1.8em}{Ausgleichung eines Vierecks}{23}{—}{245}
\IXContentsEntry{1.8em}{Gleichung zwischen den Seiten und Diagonalen eines Vierecks}{10}{—}{248}
\IXContentsText{0pt}{\textit{Briefwechsel.}}
\IXContentsText{1.8em}{Über die Wahl der Bedingungsgleichung aus den Seitenverhältnissen:}
\IXContentsEntry{.36\linewidth}{\textsc{Gauss} an \textsc{Gerling} 1824 Febr.~11}{8}{—}{249}
\IXContentsEntry{.36\linewidth}{\textsc{Gauss} an \textsc{Gerling} 1840 Jan.~19}{9}{—}{250}
\IXContentsText{0pt}{\textit{Nachlass.}}
\IXContentsEntry{1.8em}{Zur Ausgleichung der Winkel im Viereck}{18}{—}{254}
\IXContentsEntry{1.8em}{Viereck zwischen 4 Punkten 1.~2.~3.~4}{19}{—}{257}
\IXContentsEntry{1.8em}{Ausgleichung eines Polygons}{23}{—}{257}
\IXContentsEntry{1.8em}{Gewicht von Höhenbestimmungen}{21}{—}{258}
\IXContentsEntry{1.8em}{Bemerkungen}{28}{—}{259}

\IXContentsHeading{STATIONSAUSGLEICHUNGEN.}
\IXContentsText{0pt}{\textit{Nachlass.}}
\IXContentsText{1.8em}{Stationsausgleichung für Zeven aus sämmtlichen Messungen von 1824 und 1825 (ohne die}
\IXContentsEntry{3.6em}{vom 4. und 5. August)}{23}{—}{263}
\IXContentsEntry{1.8em}{Stationsausgleichung für Brillit}{22}{—}{265}
\IXContentsEntry{1.8em}{Wilsede aus sämmtlichen Messungen von 1822 und 1824}{12}{—}{267}
\IXContentsEntry{1.8em}{Ausgleichung der auf dem Windberge gemessenen Winkel}{12}{—}{271}
\IXContentsEntry{1.8em}{Beobachtungen auf Breithorn 1822}{20}{—}{274}
\IXContentsText{0pt}{\textit{Briefwechsel.}}
\IXContentsEntry{1.8em}{\rlap{Über Stationsausgleichungen:}\hspace{.32\linewidth}\textsc{Gauss} an \textsc{Gerling} 1823 Dec.~26}{10}{—}{278}
\IXContentsEntry{.36\linewidth}{\textsc{Gauss} an \textsc{Schumacher} 1827 Dec.~22}{8}{—}{281}
\IXContentsEntry{.36\linewidth}{\textsc{Gauss} an \textsc{Schumacher} 1828 Jan.~7}{9}{—}{285}
\IXContentsEntry{1.8em}{Bemerkungen}{28}{—}{287}

% Printed 526; principal PDF 533; Google PDF 537.
\par\clearpage
\noindent\makebox[0pt][l]{\fontsize{13}{16}\selectfont 526}\hfill{\small\textbf{INHALT.}}\hfill\null\par\medskip

\IXContentsHeading{ZUR NETZAUSGLEICHUNG.}
\IXContentsText{0pt}{\textit{Nachlass.}}
\IXContentsEntry{1.8em}{Anzahl der Bedingungsgleichungen in einem Dreieckssystem}{11}{Seite}{297}
\IXContentsEntry{1.8em}{Die 33 Hauptdreieckspunkte und ihre Verbindungen}{14}{—}{297}
\IXContentsEntry{1.8em}{Zusammenstellung der beobachteten Dreiecke und ihrer Widersprüche}{8}{—}{300}
\IXContentsEntry{1.8em}{Normalgleichungen, die den Winkelgleichungen entsprechen}{11}{—}{302}
\IXContentsEntry{1.8em}{Bedingungsgleichungen der zweiten Art}{19}{—}{304}
\IXContentsEntry{1.8em}{Normalgleichungen, die den Seitengleichungen entsprechen}{12}{—}{311}
\IXContentsEntry{1.8em}{Die Verbesserungen}{26}{—}{312}
\IXContentsEntry{1.8em}{Ausgleichungswerthe}{26}{—}{314}
\IXContentsEntry{1.8em}{Die Azimuthe der Seiten des sphäroidischen und des ebenen Dreieckssystems}{5}{—}{317}
\IXContentsText{0pt}{\textit{Briefwechsel.}}
\IXContentsEntry{1.8em}{\textsc{Gauss} an \textsc{Olbers} 1823 Nov.~2}{22}{—}{319}
\IXContentsEntry{1.8em}{\textsc{Gauss} an \textsc{Olbers} 1824 Juli 6}{22}{—}{320}
\IXContentsEntry{1.8em}{\textsc{Gauss} an \textsc{Olbers} 1826 Mai 14}{22}{—}{320}
\IXContentsEntry{1.8em}{\textsc{Gauss} an \textsc{Gerling} 1837 Mai 2}{22}{—}{323}
\IXContentsEntry{1.8em}{\textsc{Gauss} an \textsc{Gerling} 1838 Juni 5}{21}{—}{323}
\IXContentsEntry{1.8em}{\textsc{Gauss} an \textsc{Gerling} 1838 Nov.~14}{21}{—}{326}
\IXContentsEntry{1.8em}{Bemerkungen}{28}{—}{327}

\IXContentsHeading{DREIECKSKRANZ UM OLDENBURG.}
\IXContentsText{0pt}{\textit{Nachlass.}}
\IXContentsEntry{1.8em}{Zur Ausgleichung des Dreieckskranzes, der das Oldenburgsche umgibt}{8}{—}{331}
\IXContentsEntry{1.8em}{Bemerkungen}{28}{—}{340}

\IXContentsHeading{ZUR HANNOVERSCHEN TRIANGULATION.}
\IXContentsText{0pt}{\textit{Briefwechsel.}}
\IXContentsText{1.8em}{\textsc{Gauss} an \textsc{Schumacher}: 1816 Juli 5, 1818 Sept.~10, 1820 Mai 20, 1821 März 4, 1822 Sept.~18,}
\IXContentsEntry{3.6em}{1825 Januar, 1825 Juni 20, 1829 Jan.~14}{17}{—}{345}
\IXContentsText{1.8em}{\textsc{Gauss} an \textsc{Bessel}: 1821 Dec.~26, 1822 Nov.~15, 1823 Nov.~5, 1824 Nov.~20, 1826 März 12,}
\IXContentsEntry{3.6em}{1826 Nov.~20, 1827 April 1, 1830 April 9}{16}{—}{349}
\IXContentsEntry{1.8em}{\textsc{Gauss} an \textsc{Bohnenberger}: 1823 Nov.~16}{18}{—}{364}
\IXContentsText{1.8em}{\textsc{Gauss} an \textsc{Olbers}: 1821 Jan.~13, 1822 April 18, 1824 Juli 4, 1824 Juli 8, 1825 Febr.~19,}
\IXContentsText{3.6em}{1825 Febr.~25, 1825 Oct.~9, 1826 April 2, 1827 Jan.~14, 1827 März 1, 1830 Juni 14,}
\IXContentsEntry{3.6em}{1837 Sept.~2}{27}{—}{367}
\IXContentsText{1.8em}{\textsc{Gauss} an \textsc{Gerling}: 1821 Oct.~5, 1822 Febr.~21, 1822 Nov.~7, 1823 Juli 27, 1823 Aug.~11,}
\IXContentsEntry{3.6em}{1823 Sept.~1, 1823 Sept.~5, 1823 Oct.~3, 1827 Juli 19, 1838 Sept.~12, 1838 Nov.~14}{2}{—}{380}
\IXContentsEntry{1.8em}{Bemerkungen}{28}{—}{395}
\IXContentsText{0pt}{\textit{Veröffentlichungen.}}
\IXContentsEntry{1.8em}{Auszug aus einem Briefe des Herrn Hofrath und Ritter \textsc{Gauss}}{10}{—}{396}
\IXContentsEntry{1.8em}{Auszug. aus einem Schreiben des Herrn Hofrath \textsc{Gauss}}{13}{—}{397}
\IXContentsEntry{1.8em}{Bemerkungen}{28}{—}{400}
\IXContentsText{0pt}{\textit{Nachlass.}}
\IXContentsText{1.8em}{Plan und Anfang zum Werke über die trigonometrischen Messungen in Hannover.}
\IXContentsEntry{3.6em}{Plan des Werkes}{25}{—}{401}

% Printed 527; principal PDF 534; Google PDF 538.
\par\clearpage
\noindent\null\hfill{\small\textbf{INHALT.}}\hfill\makebox[0pt][r]{\fontsize{13}{16}\selectfont 527}\par\medskip

\IXContentsEntry{3.6em}{Einleitung}{28}{Seite}{402}
\IXContentsEntry{3.6em}{Erster Abschnitt.\quad Anordnung der Messungen im allgemeinen}{9}{—}{405}
\IXContentsText{1.8em}{Auszüge aus Berichten über die Triangulirung an das hannoversche Cabinets-Ministerium.}
\IXContentsEntry{3.6em}{Aus einem Bericht vom 7. Januar 1822 über die Arbeiten im Jahre 1821}{5}{—}{406}
\IXContentsText{3.6em}{Aus einem Bericht vom 31. Januar 1823 über die Arbeiten der Gradmessung im Jahre}
\IXContentsEntry{5.4em}{1822}{28}{—}{409}
\IXContentsText{3.6em}{Aus dem Bericht vom 16. Februar 1825 über die im Jahre 1824 ausgeführten trigono-}
\IXContentsEntry{5.4em}{metrischen Arbeiten}{23}{—}{412}
\IXContentsText{3.6em}{Aus dem Bericht vom 21. November 1827, betreffend die weitere Ausdehnung der Grad-}
\IXContentsEntry{5.4em}{messungsarbeiten}{24}{—}{413}
\IXContentsEntry{3.6em}{Aus dem Bericht vom 26. Juni 1828, die Fortsetzung der Gradmessungsarbeiten betreffend}{0}{—}{416}
\IXContentsText{3.6em}{Historischer Bericht über die von dem Hofrath \textsc{Gauss} theils ausgeführten, theils ge-}
\IXContentsEntry{5.4em}{leiteten Messungen im Königreich Hannover}{14}{—}{418}
\IXContentsText{3.6em}{Aus dem Bericht vom 5. Juli 1840 über die trigonometrischen Vermessungen im Jahre}
\IXContentsEntry{5.4em}{1839}{28}{—}{425}
\IXContentsText{3.6em}{Aus einem Bericht, December 1844, über die im Jahre 1844 ausgeführten trigono-}
\IXContentsEntry{5.4em}{metrischen Messungen}{22}{—}{425}
\IXContentsEntry{1.8em}{Geographische Coordinaten der Hauptdreieckspunkte der hannoverschen Messungen}{3}{—}{427}
\IXContentsEntry{1.8em}{Bemerkungen}{28}{—}{429}

\bigskip
\IXContentsHeading{HÖHENMESSUNGEN.}
\IXContentsText{0pt}{\textit{Veröffentlichung und Briefwechsel.}}
\IXContentsText{1.8em}{Der Refractionscoefficient aus den Höhenmessungen bei der hannoverschen Gradmessung.}
\IXContentsText{3.6em}{Beobachtete und berechnete Triangulirung im Hannoverschen, Braunschweigschen und}
\IXContentsEntry{5.4em}{Lüneburgschen}{24}{—}{437}
\IXContentsText{3.6em}{Der Refractionscoefficient aus den Höhenmessungen von 1824 abgeleitet: \textsc{Gauss} an}
\IXContentsEntry{5.4em}{\textsc{Schumacher} 1824 Nov.~28}{20}{—}{440}
\IXContentsText{3.6em}{Der Refractionscoefficient aus den Höhenmessungen bei der Gradmessung und ihrer}
\IXContentsEntry{5.4em}{Fortsetzung bis Jever: \textsc{Gauss} an \textsc{Schumacher} 1846 Dec.~27}{8}{—}{443}
\IXContentsEntry{3.6em}{Bemerkungen}{26}{—}{444}
\IXContentsText{0pt}{\textit{Nachlass.}}
\IXContentsEntry{1.8em}{Terrestrische Refraction}{25}{—}{445}
\IXContentsEntry{1.8em}{Formel zur Höhenberechnung}{22}{—}{448}
\IXContentsEntry{1.8em}{Ausgleichung der Höhen der Hauptdreieckspunkte}{15}{—}{449}
\IXContentsEntry{1.8em}{Bemerkungen}{28}{—}{454}
\IXContentsText{0pt}{\textit{Veröffentlichung.}}
\IXContentsEntry{1.8em}{Tafel für barometrisches Höhenmessen}{19}{—}{456}

\bigskip
\IXContentsHeading{HELIOTROP.}
\IXContentsText{0pt}{\textit{Veröffentlichungen.}}
\IXContentsEntry{1.8em}{Über den Heliotrop}{26}{—}{461}
\IXContentsEntry{1.8em}{Erfindung eines Heliotrops}{23}{—}{466}
\endgroup
